\documentclass{standalone}
\usepackage[T1]{fontenc}
\usepackage{mathptmx}
\usepackage{amsmath}
\usepackage{graphicx}
\usepackage{tikz}
\usetikzlibrary{calc}
\usepackage{quantikz}

\newcommand{\isw}{\raisebox{-0.5\height}[1.5pt][1.5pt]{\rotatebox{90}{\mbox{\footnotesize$\sqrt{i\text{SWAP}}$}}}}

% Single-qubit gates are SQUARE, side \gsize, and the two-qubit box is given the
% same width so all vertical edges line up.
\newlength{\gsize}
% Note: the trailing style argument is a MANDATORY brace group, not an
% optional [ ] one -- quantikz's matrix parser misreads the lookahead from
% an optional argument as a stray column/row break and stretches the gate
% to fill the whole row.  Call as \gbox{}{S^\dagger} / \gbox{name=Ca}{S^\dagger}.
%
% The box itself is a SYMMETRIC strut (gsize/2 above and below the baseline),
% giving every gate the same square footprint.  The label is measured with
% \sbox and then \raisebox by half its own (ht-dp), which moves the LABEL's
% own visual centre (not its baseline) onto the box's centre -- needed
% because "S^\dagger" (tall ascender), "H" (no descender) and "U" (no
% descender either, but a different cap height) all sit differently
% relative to their baseline, so a single fixed strut split can't centre
% all of them at once the way this per-label raise does.
\newsavebox{\glyphbox}
\newcommand{\centred}[1]{%
  \sbox{\glyphbox}{#1}%
  \raisebox{-0.5\dimexpr\ht\glyphbox-\dp\glyphbox\relax}{\usebox{\glyphbox}}}
\newcommand{\gbox}[2]{%
  \gate[style={minimum width=\gsize, inner sep=0pt, #1}]%
       {\makebox[\gsize][c]{\rule[-0.5\gsize]{0pt}{\gsize}\centred{\footnotesize\ensuremath{#2}}}}}
\newcommand{\iswbox}[1]{%
  \gate[wires=2, style={minimum width=\gsize, inner sep=0pt, #1}]%
       {\makebox[\gsize][c]{\rule[-0.5\gsize]{0pt}{\gsize}\isw}}}
% Plain two-qubit box labelled U, same square-grid convention as \gbox/\iswbox.
\newcommand{\ubox}[1]{%
  \gate[wires=2, style={minimum width=\gsize, inner sep=0pt, #1}]%
       {\makebox[\gsize][c]{\rule[-0.5\gsize]{0pt}{\gsize}\centred{\footnotesize$U$}}}}

% Small centred caption line set below each row.  Edit the wording to taste.
% Pinned to 9pt/10.5pt to match revtex4-1's native caption font, same as
% \panellabel in circuit_routed.tex -- see that file for why this is set
% explicitly rather than via \footnotesize (which is 8pt here, not 9pt).
\newcommand{\panellabel}[1]{{\fontsize{9}{10.5}\selectfont #1}}

\begin{document}
\begin{minipage}{246.0pt}
\centering
% 3.2mm is too small: S^\dagger's ascender (~6.96pt) pokes past the strut's
% height share (0.66*gsize) at that size, so its box silently grows taller
% than H's. 3.8mm gives 0.66*3.8mm=7.16pt, just enough headroom.
\setlength{\gsize}{3.8mm}%

% Row gaps below are explicit pt values, not \baselineskip multiples: array/
% tabular locally zeroes \baselineskip while building rows, so a \\[N\baselineskip]
% here silently adds no space at all.
%
% A 3-column table (symbol / "=" / circuit), all columns centred, so that
% -- column 1 centres SWAP over U,
% -- column 2 centres the two "=" signs on each other,
% -- column 3 centres the U-U-U row over the sqrt(iSWAP) block.
% Captions are inserted as \multicolumn rows spanning the whole table.
\begin{tabular}{@{}c@{\quad}c@{\quad}c@{}}
$\vcenter{\hbox{%
\begin{quantikz}[row sep={0.82cm,between origins}, column sep=0.18cm]
\qw & \qw & \swap{1} & \qw &  \qw \\
\qw & \qw & \targX{}  & \qw &  \qw
\end{quantikz}}}$
&
$=$
&
$\vcenter{\hbox{%
\begin{quantikz}[row sep={0.82cm,between origins}, column sep=0.18cm]
\qw & \ubox{} & \ubox{} & \ubox{} & \qw \\
\qw &         &         &         & \qw
\end{quantikz}}}$ \\[3pt]
\multicolumn{3}{c}{\panellabel{(a) SWAP as three applications of $U$}} \\[10pt]
$\vcenter{\hbox{%
\begin{quantikz}[row sep={0.82cm,between origins}, column sep=0.18cm]
\qw & \ubox{} & \qw \\
\qw &         & \qw
\end{quantikz}}}$
&
$=$
&
$\vcenter{\hbox{%
\begin{quantikz}[row sep={0.82cm,between origins}, column sep=0.18cm]
\qw & \gbox{}{S^{\dagger}} & \gbox{}{H} & \iswbox{} & \qw \\
\qw & \gbox{}{S^{\dagger}} & \gbox{}{H} &           & \qw
\end{quantikz}}}$ \\[3pt]
\multicolumn{3}{c}{\panellabel{\text{(b)} $U$ in terms of native gates}}
\end{tabular}

\end{minipage}
\end{document}
